\PartChapter{Detailed Theory: Signal, Automorphic, and Filtered Contact Structures}{Prime-Shell Signals and the Frey Nonlinear Transform}

\section{Prime support as bandwidth and valuation as persistence}
For
\[
 n=\prod_p p^{e_p},
\]
define the prime-shell signal
\[
 \Psi_n(p,r)
 =
 \sqrt{\log p}\,\mathbf 1_{0\le r<e_p}.
\]
Its total energy is
\[
 \norm{\Psi_n}_2^2
 =
 \sum_p e_p\log p
 =
 \boxed{\log n}.
\]
Let $P_{\mathrm{act}}$ retain only the first shell in each active prime channel. Then
\[
 \norm{P_{\mathrm{act}}\Psi_n}_2^2
 =
 \sum_{p\mid n}\log p
 =
 \boxed{\log\rad(n)}.
\]
Thus the radical is activated bandwidth, while the difference
\[
 \log n-\log\rad(n)
\]
is repeated persistence depth.

For a primitive triple, the direct-sum signal
\[
 \Psi_{a,b,c}
 =
 \Psi_a\oplus\Psi_b\oplus\Psi_c
\]
has disjoint prime support in the three sectors. The standard quality becomes the exact energy ratio
\[
 \boxed{
 q(a,b,c)
 =
 \frac{\log c}{\log\rad(abc)}
 =
 \frac{\norm{P_c\Psi_{a,b,c}}_2^2}
 {\norm{P_{\mathrm{act}}\Psi_{a,b,c}}_2^2}.
 }
\]
High quality means large output persistence relative to activated global bandwidth.

\section{Unitary transforms cannot create a radical saving}
Apply a unitary Fourier, wavelet, or polyphase transform in the shell variable. Parseval preserves total shell energy. Therefore the ratio between total depth energy and activated support energy is unchanged by a unitary shell analysis followed by exact synthesis.

\begin{proposition}[Unitary no-gain principle]
Let $U$ be unitary on the shell coordinate of every active prime channel. Then
\[
 \norm{U\Psi_n}_2=\norm{\Psi_n}_2.
\]
No inequality that depends only on these two energies becomes stronger after applying $U$.
\end{proposition}

Signal theory is still useful because it reveals which nonunitary or nonlinear operations would be required. It also separates local persistence from global support before introducing a nonlinear arithmetic carrier.

\section{The three-spike \texorpdfstring{$p$-adic}{p-adic} coherence signal}
Represent the additive triple by
\[
 \nu_{a,b,c}=\delta_0+\delta_a+\delta_c.
\]
The pairwise differences are $a$, $b$, and $c$. At a prime dividing exactly one of these integers, exactly one pair of spikes is $p$-adically close. The depth of that edge is the corresponding valuation.

The triple therefore determines a prime-indexed coherence triangle. The radical counts the active triangles. The valuations count how long a pair remains coupled through the $p$-adic refinement tree. This picture anticipates the later cross-prime logarithmic network.

\section{The Frey curve}
Attach the marked Frey curve
\[
 \boxed{
 E_{a,b,c}:y^2=x(x-a)(x+b).
 }
\]
The roots $0$, $a$, and $-b$ retain the three differences $a$, $b$, and $c$. The basic invariants are
\[
 \Delta_E=16a^2b^2c^2
\]
and
\[
 c_4(E)=16(a^2+ab+b^2).
\]
At every odd prime $p\mid abc$, the curve has multiplicative reduction. The conductor records that the channel is active, while the minimal discriminant records the full depth.

In signal language, the Frey curve is a nonlinear encoder that converts:
\[
 \begin{aligned}
 \text{prime support}&\longrightarrow\text{conductor},\\
 \text{valuation depth}&\longrightarrow\text{minimal discriminant depth},\\
 \text{additive relation}&\longrightarrow\text{elliptic geometry}.
 \end{aligned}
\]
This is precisely the kind of nonlinear transform that a unitary shell system could not provide.

\section{The automorphic eigenchannel}
Modularity associates a weight-two newform $f_E$ to the rational isogeny class of $E$. Its Hecke coefficients
\[
 a_\ell(E)
\]
form a global frequency signal. The conductor is the automorphic bandwidth. The newform, however, remembers the rational isogeny class rather than the marked integral representative.

This information loss is essential. Different integral Frey models can lie in the same rational isogeny or twist channel while having different minimal discriminant depth, N\'eron differential, and modular synthesis cost. A proof cannot use the rational eigen-signal alone.

\section{Dyadic Landen descent}
Let $u>v>0$ be coprime odd integers. Consider
\[
 T_0=(v^2,u^2-v^2,u^2)
\]
and
\[
 T_1=
 \left(
 uv,
 \left(\frac{u-v}{2}\right)^2,
 \left(\frac{u+v}{2}\right)^2
 \right).
\]
Both are primitive triples. Their radicals are equal:
\[
 \rad(T_0)=\rad(T_1)=\rad\bigl(uv(u-v)(u+v)\bigr).
\]
The second output is strictly smaller:
\[
 \left(\frac{u+v}{2}\right)^2<u^2.
\]
Hence the second quality is smaller.

The Legendre parameter of the parent is
\[
 \lambda=t^2,
 \qquad
 t=v/u.
\]
The degree-two Landen transformation gives
\[
 \boxed{
 \lambda
 \longmapsto
 \frac{4\sqrt\lambda}{(1+\sqrt\lambda)^2}.
 }
\]
This is exactly the child ratio $a_1/c_1$. The transformation preserves the automorphic channel up to the corresponding two-isogeny while lowering the integral height. It proves that the integral representative inside an isogeny class is relevant to $abc$.

\section{Exact automorphic analysis--synthesis identity}
Let
\[
 \phi_E:X_0(N_E)\longrightarrow E
\]
be an optimal modular parametrization, and let
\[
 \phi_E^*\omega_E
 =
 2\pi i\,c_E f_E(z)\,dz.
\]
Define the differential area
\[
 A(E)
 =
 \frac{i}{2}
 \int_{E(\C)}
 \omega_E\wedge\overline{\omega_E}.
\]
Pulling the area form back gives
\[
 \boxed{
 \deg(\phi_E)A(E)
 =
 4\pi^2 c_E^2
 \norm{f_E}_{\mathrm{Pet}}^2.
 }
\]
For the Frey model, scaling $x=cX$, $y=c^{3/2}Y$ gives an area factor proportional to $1/c$. Thus the height $c$ can be written as a shape-corrected synthesis gain involving modular degree, Manin normalization, and Petersson energy.

The identity does not prove a conductor-scale upper bound. It identifies the exact quantities that an integral automorphic frame theorem would have to control.

\chapter{Divisor-Polyphase Persistence and the Hecke Filter Bank}

\section{Exact-order divisor packets}
Let
\[
 n=\prod_p p^{e_p}.
\]
For every $d\ge1$, define
\[
 \boxed{
 B_d(n)
 =
 \prod_{\substack{p\mid n\\d\mid e_p}}p.
 }
\]
Since
\[
 \sum_{d\mid e}\varphi(d)=e,
\]
we have the perfect reconstruction identity
\[
 \boxed{
 n
 =
 \prod_{d\ge1}
 B_d(n)^{\varphi(d)}.
 }
\]
The first packet is
\[
 B_1(n)=\rad(n).
\]
The powerful part is
\[
 \boxed{
 \frac{n}{\rad(n)}
 =
 \prod_{d\ge2}
 B_d(n)^{\varphi(d)}.
 }
\]
This is a multiplicative polyphase decomposition indexed by exact character orders.

\section{Quality as input-output persistence imbalance}
For a primitive triple,
\[
 q-1
 =
 \frac{
 \log(c/R_c)-\log(R_aR_b)
 }{
 \log R
 }.
\]
Using the divisor packets,
\[
 \boxed{
 q-1
 =
 \frac{
 \sum_{d\ge2}\varphi(d)\log B_d(c)
 -
 \log(R_aR_b)
 }{\log R}.
 }
\]
A high-quality triple has more repeated-depth persistence in the output than the input radicals supply as activated bandwidth.

Strong $abc$ is equivalent to
\[
 \prod_{d\ge2}
 B_d(c)^{\varphi(d)}
 \le
 C_\eps R_aR_bR^\eps.
\]
The decomposition makes the target transparent, but not easier by itself. A global channel must explain why the exact-order packets cannot all be large simultaneously.

\section{Local Tate-depth interpretation}
At an odd prime $p^e\parallel c$, the Frey discriminant has local depth $2e$. The divisors of $e$ organize the levels at which prime-power coefficient systems can become insensitive to the local extension. The exact-order divisor packet therefore has a local representation-theoretic interpretation, although local disappearance does not automatically produce a global low-level eigenform.

\section{Atkin--Lehner Walsh packets}
Let
\[
 W_1,\ldots,W_r
\]
be commuting Atkin--Lehner involutions. For a sign vector
\[
 \varepsilon\in\{\pm1\}^r,
\]
define
\[
 A_\varepsilon
 =
 \prod_{j=1}^{r}(I+\varepsilon_jW_j).
\]
Then
\[
 \boxed{
 A_\varepsilon^2=2^rA_\varepsilon,
 }
\]
\[
 A_\varepsilon A_{\varepsilon'}=0
 \quad(\varepsilon\neq\varepsilon'),
\]
and
\[
 \boxed{
 \sum_\varepsilon A_\varepsilon=2^rI.
 }
\]
The unnormalized analysis packets are integral. The common synthesis denominator is $2^r$, which is subpolynomial in a squarefree conductor. Local sign separation is therefore not the large persistence carrier.

The automorphic interpretation uses the standard Atkin--Lehner and modular-newform framework. The integral separation problem is related to congruence modules and modular degrees~\cite{Wiles,TaylorWiles,BCDT,Agashe,CesnaviciusManin}.

\section{The same-sign separation problem}
After fixing the local signs, several eigenforms can remain congruent inside the same packet. Let $\Lambda$ be an integral Hecke lattice and let $e_f$ be the rational projector onto the Frey eigenline. Define
\[
 \eta_f
 =
 \min\{m\ge1:m e_f(\Lambda)\subseteq\Lambda\}.
\]
This is the integral eigenchannel separation denominator. It is an arithmetic analogue of leakage in a filter bank. The difficulty is not rational spectral isolation. It is isolation inside the integral lattice.

\chapter{Integral Projectors, Cotangent Multiplicity, and Wiles Defect}

\section{Projector denominator as a congruence exponent}
Let $V=U\oplus W$ be a rational decomposition and let $\Lambda\subset V$ be an integral lattice. Let $e_U$ be the rational projector onto $U$. Define
\[
 \Lambda_U=\Lambda\cap U,
 \qquad
 \Lambda_W=\Lambda\cap W,
\]
and
\[
 C_U=\Lambda/(\Lambda_U+\Lambda_W).
\]

\begin{theorem}[Integral projector theorem]
The least positive integer $\eta_U$ satisfying
\[
 \eta_U e_U(\Lambda)\subseteq\Lambda
\]
is the exponent of the finite quotient $C_U$.
\end{theorem}

\begin{proof}
For $x=u+w\in\Lambda$, the condition $m e_U(x)=mu\in\Lambda_U$ is equivalent to
\[
 mx=mu+mw\in\Lambda_U+\Lambda_W.
\]
Thus $m$ clears every projector denominator exactly when it annihilates every class in $C_U$.
\end{proof}

For a primitive rank-one vector $v\in\Z^n$, the orthogonal projector denominator is
\[
 \langle v,v\rangle,
\]
and the congruence quotient is cyclic of that order.

\section{Why one cyclic quotient is not enough for split persistence}
The split persistence module
\[
 \mathcal P_c
 \cong
 \bigoplus_{p^e\parallel c}
 (\Z/p\Z)^{e-1}
\]
has large $p$-rank. The $p$-torsion of a cyclic target has rank at most one. Therefore every map from $\mathcal P_c$ to a cyclic group has a forced kernel of order at least
\[
 \prod_{p^e\parallel c}p^{\max(e-2,0)}.
\]
This is a true obstruction to the split source. It motivates the Hecke cotangent module
\[
 \mathfrak p/\mathfrak p^2,
\]
which can retain multiple independent congruence directions before cyclic compression.

\section{A toy Hecke order}
Consider
\[
 \mathbb T(\mathbf m)
 =
 \{(x_0,x_1,\ldots,x_r)\in\Z^{r+1}:x_i\equiv x_0\pmod{m_i}\}.
\]
For the target coordinate $x_0$, the cotangent module is
\[
 \boxed{
 \mathfrak p/\mathfrak p^2
 \cong
 \bigoplus_{i=1}^{r}\Z/m_i\Z.
 }
\]
The cyclic congruence quotient is
\[
 \boxed{
 \Z/\lcm(m_1,\ldots,m_r)\Z.
 }
\]
The quotient of cotangent order by congruence order is
\[
 \frac{\prod_i m_i}{\lcm(m_i)}.
\]
This is a toy Wiles-defect factor. It measures the multiplicity discarded by cyclic compression.

The comparison between cotangent and congruence data, including non-complete-intersection defects and local--global formulas, is developed in the Wiles-defect literature~\cite{BKM1,BKM2}.

\section{Complete-intersection control case}
Let
\[
 \mathbb T=\Z[X]/(F)
\]
and evaluate at a simple root $a_0$. Then the congruence ideal and cotangent module both have order
\[
 |F'(a_0)|.
\]
The defect is one. Extra multiplicity is therefore tied to non-complete-intersection structure rather than being a normalization artifact.

\section{Two prime indices}
The persistence is indexed by conductor primes $q\mid abc$. Congruence depth is indexed by coefficient primes $\ell$. The natural local signal is a matrix
\[
 D_f(\ell,q).
\]
At a split Tate prime $q\neq\ell$, scalarity modulo $\ell^n$ forces
\[
 n\le\min\{v_\ell(q-1),v_\ell(v_q(q_E))\}.
\]
These off-diagonal channels carry tame congruence information. They do not naturally reproduce the full base-$q$ factor $q^{e-1}$. This mismatch motivates the later filtered contact correction and the diagonal coefficient-prime channel.

\chapter{Filtered Cusp Contact and Witt Carry Structure}

\section{The three-cusp section}
View a primitive triple as the integral section
\[
 \sigma_{a,b,c}:
 \Spec\Z\longrightarrow\mathbb P^1_\Z,
 \qquad
 \sigma=[a:c].
\]
The pullbacks of the boundary divisors $0$, $1$, and $\infty$ are
\[
 \operatorname{div}(a),
 \qquad
 \operatorname{div}(b),
 \qquad
 \operatorname{div}(c).
\]
For the output cusp, let
\[
 C=\operatorname{div}(c),
 \qquad
 C_{\mathrm{red}}=\operatorname{div}(R_c).
\]
Define the two-term contact complex
\[
 \boxed{
 \mathbf J_c
 =
 [\mathcal O(-C)\longrightarrow\mathcal O(-C_{\mathrm{red}})].
 }
\]
On global sections,
\[
 \mathbf J_c
 =
 [c\Z\hookrightarrow R_c\Z]
 \simeq
 [\Z\xrightarrow{c/R_c}\Z].
\]

\begin{theorem}[Filtered contact quotient]
\[
 \boxed{
 H^0(\mathbf J_c)
 =
 R_c\Z/c\Z
 \cong
 \Z/(c/R_c)\Z.
 }
\]
\end{theorem}

The output persistence is cyclic before grading.

\section{Associated graded and carry extensions}
At $p^e\parallel c$,
\[
 \Jcal_{p,e}
 =
 p\Z_p/p^e\Z_p
 \cong
 \Z/p^{e-1}\Z.
\]
Its filtration
\[
 p^j\Z_p/p^e\Z_p
\]
has successive quotients $\F_p$. Thus
\[
 \gr\Jcal_{p,e}\cong(\F_p)^{e-1}.
\]
The exact sequences
\[
 0\longrightarrow
 \Z/p^{j-1}\Z
 \xrightarrow{\times p}
 \Z/p^j\Z
 \longrightarrow
 \F_p
 \longrightarrow0
\]
are nonsplit. These carry extensions are precisely what the split persistence packet forgot.

\begin{corollary}[Correction of the cyclic-rank barrier]
The $p$-rank obstruction applies to the associated graded source. It does not obstruct a filtered map from $\Z/p^{e-1}\Z$ to a cyclic target of matching depth.
\end{corollary}

\section{Exact cyclic transducer}
Let
\[
 C(m)=[\Z\xrightarrow{m}\Z]
\]
and
\[
 C(\eta)=[\Z\xrightarrow{\eta}\Z].
\]
Set $g=\gcd(m,\eta)$. Multiplication by $m/g$ in degree $-1$ and by $\eta/g$ in degree $0$ defines a chain map. On cohomology,
\[
 \boxed{
 |\ker|=m/g,
 \qquad
 |\im|=g,
 \qquad
 |\coker|=\eta/g.
 }
\]
The maximal transferable packet is the common prime-power depth.

For the output source $m=c/R_c$, the corrected algebraic defect against a cyclic target of order $\eta_f$ is
\[
 \boxed{
 \delta_{\mathrm{alg}}(c,f)
 =
 \frac{c/R_c}{\gcd(c/R_c,\eta_f)}.
 }
\]
Prime by prime,
\[
 \delta_{\mathrm{alg}}(c,f)
 =
 \prod_{p^e\parallel c}
 p^{\max(e-1-v_p(\eta_f),0)}.
\]
The problem is missing depth, not missing rank.

\section{Witt realization}
The local packet has a natural additive identification
\[
 \boxed{
 \Jcal_{p,e}
 \cong
 W_{e-1}(\F_p).
 }
\]
The Verschiebung filtration matches multiplication by successive powers of $p$. This is the typed local object underlying the carry chain.

At odd $p$, the $p$-adic logarithm identifies
\[
 (1+p\Z_p)/(1+p^e\Z_p)
 \cong
 p\Z_p/p^e\Z_p.
\]
At $2$, the additive contact packet remains uniform, while the multiplicative coordinate begins on $1+4\Z_2$ and therefore carries the familiar weighted dyadic shift.

\section{The filtered local-global target}
The intended transducer is a filtered chain
\[
 \mathbf J_c
 \longrightarrow
 \bigoplus_{p\mid c}\mathbf L_p^{\mathrm{st}}
 \longrightarrow
 \mathbf L_{R^{\mathrm{glob}}}
 \xrightarrow{R=\mathbb T}
 \mathbf L_{\mathbb T_f}
 \longrightarrow
 \mathbf C_f.
\]
It must preserve the Bockstein carry maps. The local contact packet exists canonically. The unproved step is its transport through global deformation and Hecke complexes with small kernel and controlled archimedean size.

\begin{theorem}[Conditional filtered-transfer criterion]
Suppose a canonical map
\[
 H^0(\Theta_f):\Z/(c/R_c)\Z\longrightarrow H^0(\mathbf C_f)
\]
satisfies, for every $\eps>0$,
\[
 |\ker H^0(\Theta_f)|\le C_\eps R^\eps
\]
and
\[
 |H^0(\mathbf C_f)|\le C'_\eps R_aR_bR^\eps.
\]
Then the strong $abc$ conjecture follows.
\end{theorem}

\begin{proof}
The image has order
\[
 \frac{c/R_c}{|\ker|}
\]
and is bounded by the target order. Therefore
\[
 \frac{c/R_c}{C_\eps R^\eps}
 \le
 C'_\eps R_aR_bR^\eps.
\]
Multiplying by $R_c$ and replacing $\eps$ by $\eps/2$ gives the strong bound.
\end{proof}
